\documentclass[a4,11pt]{aleph-notas}
% Se puede ver la documentación aquí: 
% https://github.com/alephsub0/LaTeX_aleph-notas

% -- Paquetes adicionales 
\usepackage{enumitem}
\usepackage{aleph-comandos}
\usepackage{booktabs}
\usepackage{array}
\newcolumntype{L}[1]{>{\raggedright\arraybackslash}p{#1}}
\hypersetup{
    urlcolor=blue,
    linkcolor=blue,
}

% -- Datos 
\institucion{Escuela de Ciencias Físicas y Matemática}
\carrera{Ciencia de Datos}
\asignatura{Aprendizaje Automático Inicial}
\tema{ApS: Aprendizaje supervisado}
\autor{Andrés Merino}
\fecha{Periodo 2025-2}

\logouno[0.14\textwidth]{Logos/logoPUCE_04_ac}
\definecolor{colortext}{HTML}{0030A1}
\definecolor{colordef}{HTML}{0030A1}
\fuente{montserrat}


% -- Comandos adicionales
\setlist[enumerate]{label=\roman*.}


\begin{document}

\encabezado

\section{Resultado de Aprendizaje}

RdA 3. Resolver problemas prácticos mediante el uso de modelos de aprendizaje automático, ajustándolos para la mejora de su rendimiento y precisión.

\begin{itemize}[leftmargin=*]
\item 
    En esta actividad se evalúa si el estudiante \textit{Criterio 3.2: Aplica modelos de aprendizaje supervisado en escenarios del mundo real, ajustando los modelos para maximizar su precisión y eficiencia mediante técnicas de ajuste de hiperparámetros y regularización.}
\end{itemize}


%%%%%%%%%%%%%%%%%%%%%%%%%%%%%%%%%%%%%%%%
\section{Tema}
%%%%%%%%%%%%%%%%%%%%%%%%%%%%%%%%%%%%%%%%

\begin{center}
    \large\bf Predicción de puntos críticos de libadores alrededor de la PUCE
\end{center}


\section{Ficha del proyecto}
\begin{itemize}
    \item \textbf{Problema:} Convivencia en vía pública -- predicción de puntos críticos de libadores en calles circundantes a la PUCE.
    \item \textbf{Comunidad:} Moradores y juntas barriales de La Mariscal y La Floresta (calles aledañas a la PUCE).
    \item \textbf{Ámbito geográfico:} Radio de \(\sim\)800 m alrededor del campus.
\end{itemize}

\section{Objetivo general}
Construir y validar un modelo de {clasificación} que estime la {probabilidad de presencia de libadores} por {segmento de calle \(\times\) franja horaria} y entregar {insumos accionables} (mapa, plan de patrullaje/convivencia y piezas de comunicación), co-diseñados con las {juntas barriales y moradores}.


\subsection*{Objetivos específicos (alineados al RdA)}
\begin{itemize}
    \item Modelar el fenómeno con {Árboles de Decisión}, {SVM} (con ajuste de hiperparámetros) y {k-NN}, comparando su desempeño.
    \item Ajustar y optimizar los modelos (búsqueda de \(C\), \(\gamma\), \emph{kernel}; elección de \(k\) y métrica de distancia; poda/profundidad de árboles) maximizando {PR-AUC} y {Recall} de la clase positiva.
    \item Segmentar {microzonas} con {k-means} para orientar estrategias diferenciadas (control, mediación, campañas).
    \item Transferir los resultados mediante {dashboard/mapa}, {plan de patrullaje por franja horaria} y {kit de comunicación}.
\end{itemize}

\section{Itinerario A--E (Aprendizaje--Servicio)}

\subsection*{A. Motivación (Semana 1)}
\begin{itemize}
    \item Conversatorio con vecino (testimonio).
    \item Mini--\emph{demo}: Árbol de decisión \emph{baseline} sobre datos sintéticos para discutir riesgo vs.\ probabilidad y costos de error.
    \item {Productos:} Protocolo de seguridad en campo y acuerdo de convivencia.
    \item {Reflexión:} Diario: «¿Por qué importa este problema?».
\end{itemize}

\subsection*{B. Diagnóstico (Semana 2)}
\begin{itemize}
    \item {Prediagnóstico docente:} mapeo de capas base (OSM, alumbrado, ciclovías, zonas comerciales) y lineamientos con socios.
    \item {Diagnóstico estudiantil:} 
    \begin{itemize}
        \item {Checklist de observación} por franjas horarias: iluminación, actividad comercial, flujo peatonal, presencia policial, residuos.
        \item {Levantamiento de puntos}: conteo anónimo de grupos con consumo visible (\(\geq\)2 personas, \(\geq\)10--15 min).
        \item {Encuesta breve} a residentes/comerciantes y a libadores.
    \end{itemize}
    \item {Reflexión:} Diario «Lo que observé».
\end{itemize}

\subsection*{C. Diseño y Planificación (Semana 3)}
\begin{itemize}
    \item {Variable objetivo:} Clase {1 (alto)} si en el {segmento \(\times\) franja} se observa \(\geq\)1 evento con \(\geq\)2 personas consumiendo por \(\geq\)10--15 min {o} hay \(\geq N\) reportes en bitácora; clase {0 (no alto)} en caso contrario.
    \item {Características (ejemplos):}
    \begin{itemize}
        \item Infraestructura: iluminación (binaria/escala), presencia de ciclovía, ancho de acera, cercanía a parques/plazas, pendiente.
        \item Entorno comercial: densidad/proximidad de bares, licoreras y tiendas 24h.
        \item Dinámica temporal: día de semana, franja horaria, eventos especiales.
        \item Seguridad: proximidad a UPC/patrullaje.
        \item Clima: lluvia/no lluvia.
    \end{itemize}
    \item {Partición y validación:} \emph{Blocked/Time series CV} por semanas; estratificación por segmentos; separación espacial mínima para evitar \emph{leakage}.
    \item {Plan de ajuste:}
    \begin{itemize}
        \item {SVM:} búsqueda de \(C\), \(\gamma\), \emph{kernel} (RBF/lineal), \emph{class weights}.
        \item {k-NN:} \(k\) y métrica (euclídea/manhattan), normalización.
        \item {Árbol:} profundidad, \texttt{min\_samples\_split}, criterio Gini/Entropía.
    \end{itemize}
    \item {Reflexión:} Diario ``Decisiones de diseño''.
\end{itemize}

\subsection*{D. Ejecución (Semanas 4--6)}
\begin{itemize}
    \item Limpieza y consolidación del \emph{dataset}; ingeniería de características (ventanas temporales, distancias, densidades).
    \item Entrenamiento y {tuning} de SVM, k-NN y Árbol; selección de modelo con {PR-AUC}, {F1-macro} y {Recall positivo}; análisis de {umbrales} con matriz de costos (FN \(>\) FP).
    \item Explicabilidad: reglas del árbol e importancia de variables; \emph{model card} con supuestos y límites.
    \item Prototipos:
    \begin{itemize}
        \item Mapa/heatmap por franja (Folium/Kepler/QGIS).
        \item Plan de patrullaje sugerido (bloques de 2h, rutas y puntos).
        \item Piezas de comunicación (afiche/infografía).
    \end{itemize}
\end{itemize}

\subsection*{E. Cierre (Semana 7)}
\begin{itemize}
    \item {Asamblea vecinal:} demo del mapa, discusión de umbrales y priorización conjunta de acciones.
    \item Comunicación pública: \emph{one-pager} ejecutivo e infografías.
    \item Sistematización: informe final, \emph{model card}, \emph{datasheet} del conjunto de datos .
    \item {Reflexión:} Diario «Impacto y límites».
\end{itemize}

\section{Experiencias de aprendizaje}
\begin{enumerate}[label=\textbf{\arabic*.}]
    \item \textbf{Diagnóstico en campo} 
    \begin{itemize}
        \item \textbf{Objetivo:} Levantar datos por \textit{segmento × franja}.
        \item \textbf{Productos:} Dataset crudo georreferenciado (observaciones y encuestas) + bitácora de salida.
    \end{itemize}

    \item \textbf{Modelado y ajuste}
    \begin{itemize}
        \item \textbf{Objetivo:} Entrenar y comparar Árbol de Decisión, k-NN y SVM; ajustar hiperparámetros y umbral con matriz de costos (FN > FP).
        \item \textbf{Productos:} Tabla comparativa de métricas (PR-AUC, F1-macro, Recall positivo), curvas PR y análisis de umbral; \emph{notebook} reproducible.
    \end{itemize}

    \item \textbf{Traducción a intervención y socialización} 
    \begin{itemize}
        \item \textbf{Objetivo:} Segmentar microzonas con k-means y co-definir acciones (patrullaje, mediación, comunicación) con moradores.
        \item \textbf{Productos:} Heatmap por franja; microzonas k-means; plan de patrullaje por bloques de 2h; \emph{one-pager} y afiche barrial.
    \end{itemize}
\end{enumerate}


\section{Cronograma y logística (7 semanas)}
\begin{center}
    \renewcommand{\arraystretch}{1.5}
\begin{tabular}{@{}cL{2cm}L{6.5cm}L{5cm}@{}}
\toprule
Semana & Dónde & Actividades& Productos \\
\midrule
1 & Aula (1) & Motivación, ética, \emph{baseline}; equipos y roles & Protocolo de seguridad; Diario 1 \\
2 & Aula (1) {Campo (2)} & Checklist; recorridos & Datos de observación; Diario 2 \\
3 & Aula (1) & Variable, \emph{features}, validación y plan de ajuste & Plan validado; Diario 3 \\
4 & Aula (1) & Limpieza/EDA; primer entrenamiento & Bitácora de experimentos \\
5 & Aula (1) & {Tuning} SVM/k-NN/Árbol; análisis de umbral & Resultados comparados \\
6 & Aula (1) {Campo (1)} & k-means; prueba de campo del mapa; prototipos & Mapa preliminar; Diarios 4--5 \\
7 & Aula (1) & Socialización con socios; cierre y sistematización & Informe final, \emph{one-pager}, Diario 6 \\
\bottomrule
\end{tabular}
\end{center}

\section{Entregables finales}
\begin{itemize}
    \item Informe técnico (problema, datos, metodología, métricas, comparación, umbral, límites).
    \item Mapa/tablero de probabilidad por franja y microzonas k-means.
    \item {Kit de comunicación vecinal}  para difundir horarios y puntos sensibles.
    \item Plan de patrullaje sugerido.
    \item Diario reflexivo completo.
\end{itemize}

\section{Evaluación}

\subsection*{Rúbrica (niveles 4--1)}
\paragraph{1) Producto técnico (50\%)}~
\begin{itemize}
    \item {4 Excelente:} \emph{Pipeline} reproducible; compara SVM, k-NN y Árbol con validación bloqueada; justifica el \emph{modelo elegido} con {PR-AUC} \(\geq 0.75\) y {Recall positivo} \(\geq 0.80\) (referenciales); ajuste de umbral con matriz de costos; análisis de errores y \emph{ablation} de \emph{features}.
    \item {3 Bueno:} Comparación completa con \emph{tuning} básico; PR-AUC \(\geq 0.65\) y Recall \(\geq 0.70\); reporte claro sin \emph{ablation}.
    \item {2 Básico:} Uno o dos modelos; validación insuficiente; métricas sin análisis de umbral.
    \item {1 Insuficiente:} Modelo único sin validación; métricas inconsistentes o irrelevantes.
\end{itemize}

\paragraph{2) Impacto y pertinencia (10\%)}~
\begin{itemize}
    \item {4:} Recomendaciones accionables, priorizadas y validadas con socios; mapa claro por franjas; plan de patrullaje factible.
    \item {3:} Entregables útiles con validación parcial.
    \item {2:} Entregables poco conectados a necesidades reales.
    \item {1:} Sin utilidad práctica evidente.
\end{itemize}

\paragraph{3) Proceso y registro (10\%)}~
\begin{itemize}
    \item {4:} Bitácora completa; datos limpios con \emph{data dictionary}; cumplimiento estricto de ética, seguridad y anonimización.
    \item {3:} Registro continuo con pequeñas lagunas; ética atendida.
    \item {2:} Registros incompletos; control ético parcial.
    \item {1:} Ausencia de registros; incumplimientos.
\end{itemize}

\paragraph{4) Reflexión crítica (10\%)}~
\begin{itemize}
    \item {4:} Diarios profundos (6 entradas) con aprendizaje técnico--ético y conexión comunitaria.
    \item {3:} Diarios constantes con análisis suficiente.
    \item {2:} Diarios descriptivos sin análisis.
    \item {1:} Diarios ausentes o mínimos.
\end{itemize}

\paragraph{5) Presentación y comunicación (10\%)}~
\begin{itemize}
    \item {4:} Exposición clara para audiencia no técnica; \emph{one-pager} ejecutivo, afiches y demo funcional.
    \item {3:} Presentación adecuada con materiales completos.
    \item {2:} Presentación poco clara; materiales limitados.
    \item {1:} Presentación deficiente.
\end{itemize}




\end{document}